\documentclass[10pt]{article}\usepackage[margin=.7in]{geometry}\usepackage[T1]{fontenc}\setlength{\parindent}{0pt}\setlength{\parskip}{8pt}\sloppy\begin{document}Polynomial remainder cancellation trace invariant

T-1190; 2026-10-31. Producer proof packet; independent source-fidelity review pending.

A finite trace consists of exponent/coefficient pairs (j,c). Each step replaces r by r-c X\textasciicircum{}j b. The trace quotient is the sum of these c X\textasciicircum{}j terms. The invariant states r\_initial=b*trace\_quotient+r\_final. Divisibility by b is unchanged at every step and over the full trace. Hence zero final remainder implies b divides the initial polynomial, with the trace quotient as an explicit witness. Specialization to f-v\textasciicircum{}2 gives the divisibility clause needed for Mumford validity.

The statement admits arbitrary finite traces over a field, without assumptions on chosen exponents or coefficients. It proves soundness of a zero result, not that a particular trace terminates, computes the unique reduced remainder, or returns zero whenever the input is divisible. The native quotient array uses assignment q[j]=c, while the mathematical witness sums contributions. Their equality needs a separate distinct-exponent argument. prem returns only the native remainder, so its soundness can use this witness without asserting that the overwritten quotient is correct.

Seven main declarations and eight controls cover complete and partial division of X\textasciicircum{}2-1 by X-1, the quotient X+1, empty traces, repeated quotient slots, divisibility without a zero nonterminal remainder, and a zero divisor. Repeated exponent zero with two coefficients one produces trace quotient two, not the final assigned coefficient one. This control prevents claiming native quotient correspondence from the trace invariant alone.

Native source pdiv subtracts c*b[k] at each index j+k then trims the result; prem selects the remainder. The retained exact source excerpt covers trim and polynomial arithmetic through pexact. To close the native validity bridge, prove that each bounded UInt64/list update decodes to this polynomial cancellation, that trimming preserves that decoding, and that the accepted empty vector decodes to zero. Termination and canonical residue bounds remain separate source obligations. Existing upstream trim-preservation proofs may be reused with an exact representation map.

Sources: inspected three\_leads.rs lines795-889; pinned Mathlib Algebra/Ring/Divisibility/Basic.lean dvd\_sub\_left; prior stored-validity and canonical arithmetic packets. No new compiled dependencies. Three final native checks passed with unchanged input hashes. All 15 separate readbacks use only propext, Classical.choice and Quot.sound. Failed attempts (two invariant rewrites and one control type-inference/cast issue) are retained; outcomes are in status.json. Existing compiled upstream dependencies retain their trust boundary. No scientific execution, timing observations or cost ratios.
\end{document}